\chapter{Технологический раздел}

\section{Выбор средств реализации}

Наиболее важным внешним компонентом в разрабатываемом методе является морфологический анализатор. Предпочтительным для использования в разработанном методе является анализатор, основанный на словарном подходе, поскольку анализаторы данного типа не требуют обработки контекста. Среди анализаторов основанных на словарном подходе можно выделить следующие:
\begin{itemize}
	\item MyStem~\cite{segalovich2003fast}~---~морфологический анализатор разработанный И.~Сегаловичем, реализованный на языке Си и доступный в качестве утилиты командной строки и подключаемой библиотеки для C/C++ и Python;
	\item Stemka~\cite{StemkaWebsite}~---~морфологический анализатор разработанный А.~Коваленко, реализованный на языке С++, доступный в качестве утилиты командной строки и подключаемой библиотеки для C++;
	\item xMorphy~\cite{Sapin_2023}~---~морфологический анализатор реализованный на языке C++ и доступный в качестве утилиты командной строки и подключаемой библиотеки для C++;
	\item PyMorphy~\cite{Pymorphy2Internals}~---~морфологический анализатор, реализованный на языке Python и доступный в качестве подключаемой библиотеки для данного языка.
\end{itemize}

Для сравнения приведенных морфологических анализаторов были сформулированы следующие критерии:
\begin{itemize}
	\item доступность в качестве подключаемой библиотеки;
	\item возможность получения необходимых морфологических характеристик для каждой из рассматриваемых частей речи;
	\item возможность получения нескольких вариантов морфологических характеристик.
\end{itemize}

Сравнение перечисленных морфологических анализаторов по сформулированным критериям приведено в таблице~\ref{tab:compare_morph}.

\begin{table}[H]
	\centering
	\caption{Сравнение систем синтаксической разметки}
	\begin{tabular}{|p{4cm}|p{2cm}|p{2cm}|p{2cm}|p{2cm}|}
		\hline
		Критерий & MyStem & Stemka & xMorphy & PyMorphy\\\hline
		Использование в качестве подключаемой библиотеки & Да & Да  & Да & Да\\\hline
		Получение необходимых морфологических характеристик & Да & Нет & Да & Да \\\hline
		Получение нескольких возможных вариантов морфологических характеристик & Нет & Да & Да & Да \\\hline
	\end{tabular}
	\label{tab:compare_morph}
\end{table}

На основании приведенного сравнения анализаторы xMorphy и PyMorphy являются наиболее подходящими для использования в разработанном методе, как наиболее полно удовлетворяющие требованиям.

В качестве языка реализации метода был выбран Python, поскольку данный язык программирования поддерживает реализации необходимых структуры данных, предоставляет возможность использования web-фреймворков для реализации серверной части web-приложения и совместим с выбранным в качестве одного из наиболее подходящих для использования в реализуемом методе морфологическим анализатором PyMorphy.

\section{Пользовательский интерфейс}

\section{Реализация ключевых этапов метода}

Реализация алгоритма получения морфологических характеристик словоформ и определения типов знаков препинания приведена в листинге~\ref{lst:morph}.

\begin{breakablecode}[label={lst:morph}]{Реализация алгоритма получения морфологических характеристик словоформ и типов знаков препинания}
def get_features(word):
	punct_feats = _get_punctuation_features(word)
	if punct_feats is not None:
		return punct_feats
		
	digit_feats = _get_digit_features(word)
	if digit_feats is not None:
		return digit_feats
	
	parses = morph.parse(word)
	filtered = [p for p in parses if p.score >= MIN_PARSE_SCORE]
	if not filtered:
		filtered = parses[:1]
	
	results = []
	for p in filtered:
		feats = {'pos': POS_MAP.get(p.tag.POS, 'X')}
		g = p.tag.grammemes
		
		gender = _extract_gender(g)
		if gender is not None:
			feats['gender'] = gender
		number = _extract_number(g)
		if number is not None:
			feats['number'] = number
		case = _extract_case(g)
		if case is not None:
			feats['case'] = case
		person = _extract_person(g)
		if person is not None:
			feats['person'] = person
		tense = _extract_tense(g)
		if tense is not None:
			feats['tense'] = tense
		
		verb_feats = _extract_verb_features(g, p.tag.POS)
		feats.update(verb_feats)
		
		results.append(feats)
		
	if len(results) == 1 and results[0].get("pos") == "X":
		core = word.replace("ё", "е").replace("Ё", "Е")
		if len(core) >= 2 and core[0].isupper() and core.replace("-", "").isalpha():
			return [{"pos": "N", "number": "sg", "case": "nomn"}]
	return results
\end{breakablecode}

В листинге~\ref{lst:cyk} приведена реализация алгоритма построения таблицы составляющих. Данная реализация использует функции для вывода всех нетерминало по унарным правилам и проверки согласования, реализации которых представлены в листингах~\ref{lst:unary} и~\ref{lst:agree} соответственно.

\begin{breakablecode}[label={lst:cyk}]{Реализация алгоритма построения таблицы составляющих}
def build_cyk_table(token_feature_pairs, unary_index, binary_index):
	n = len(token_feature_pairs)
	tokens = [pair[0] for pair in token_feature_pairs]
	dp = [[dict() for _ in range(n+1)] for _ in range(n+1)]
	
	for j in range(1, n+1):
		token, features_list = token_feature_pairs[j-1]
		for feat in features_list:
			pos = feat['pos']
			feat_set = frozenset(feat.items())
			dp[j-1][j].setdefault(pos, set()).add(feat_set)
			for lhs in unary_index.get(pos, []):
				dp[j-1][j].setdefault(lhs, set()).add(feat_set)
		
	while _unary_closure_round(dp, unary_index, tokens, n):
		pass
		
	for length in range(2, n + 1):
		for i in range(0, n - length + 1):
			j = i + length
			for k in range(i + 1, j):
				left_cell = dp[i][k]
				right_cell = dp[k][j]
				for symA, featsA_set in left_cell.items():
					for symB, featsB_set in right_cell.items():
						for lhs in binary_index.get((symA, symB), []):
							for featA in featsA_set:
								for featB in featsB_set:
									if agreement_check(
									lhs, symA, featA, symB, featB,
									tokens=tokens,
									span_left=(i, k),
									span_right=(k, j),
									):
										new_feat = merge_features(lhs, symA, featA, symB, featB)
										dp[i][j].setdefault(lhs, set()).add(new_feat)
		while _unary_closure_round(dp, unary_index, tokens, n):
			pass
	return dp
\end{breakablecode}